\documentclass[../../../include/open-logic-section]{subfiles}

\begin{document}

\olfileid{his}{set}{mythology}

\olsection{تر تاريخ زيات اسطوري بيان؟}

چې له يوې پېړۍ څخه د زيات وخت له واټنه دې پېښو ته بېرته ګورو، بايد پام وکړو چې دا پرېکړې بې له پوښتنې ونۀ منو۔ نتيجې بې‌شکه نوې، په زړه پورې او حيرانوونکې وې۔ خو په رښتيا څومره ټکان ورکوونکې وې؟ او آيا په رښتيا يې ښودله چې پر هندسي وجدان بايد تکيه ونۀ کړو؟

د ټکان په مسئله کښې، \citet{Gouvea2011} څرګندوي چې ډېډېکنډ ته د کانتور مشهور عبارت، «\emph{je le vois, mais je ne le crois
pas}»، تر ډېره له خپل سياق څخه بهر اخيستل کېږي۔ د همدغه سياق يو څه زياته برخه دا ده (له \citeauthor{Gouvea2011} څخه نقل شوې ده):
\begin{quote}
که ستاسو پر مهربانۍ او زغم دومره غوښتنې کوم، نو د دې موضوع لپاره زما توده مينه وبخښئ؛ هغه خبرې چې په دې وروستيو کښې مې درلېږلې دي، ان زما لپاره هم دومره ناڅاپي او دومره نوې دي چې تر هغو د زړه سکون نۀ شم موندلے، څو له تاسو، زما قدرمن ملګري، د هغو د سموالي په اړه پرېکړه ترلاسه نۀ کړم۔ تر څو چې تاسو راسره موافق شوي نۀ يئ، يوازې دومره ويلے شم: \emph{je le vois, mais je ne le crois pas.} [وينم يې، خو باور پرې نۀ کوم۔]
\end{quote}
کانتور پوهېده چې نتيجه يې «دومره ناڅاپي، دومره نوې» وه۔ خو دا شکمنه ده چې هغه کله هم خپله نتيجه \emph{د باور وړ نۀ} ګڼلې وي۔ لکه چې \citeauthor{Gouvea2011} څرګندوي، هغه يوازې له ډېډېکنډ څخه غوښتل چې وړاندې کړے ثبوت يې وګوري۔

د هندسي وجدان په مسئله کښې: پېانو خپل ځای ډکوونکے منحني خط له هېڅ انځور پرته خپور کړ۔ خو چې هېلبرټ خپل منحني خط خپور کړ، نو خپله موخه يې څرګنده کړه: لوستونکو ته به د پېانو د نتيجې د پوهېدو يوه روښانه لاره ورکړي، که هغوي «له لاندې هندسي وجدان څخه مرسته واخلي»؛ ورپسې يې د \emph{انځورونو} يوه لړۍ وړاندې کړه، عين لکه هغه چې په \olref[his][set][pathology]{sec} کښې ورکړل شوي دي۔

په عمومي توګه: که څه هم په خپرو شويو ثبوتونو کښې د انځورونو کارول يو څه له رواجه لوېدلي، خو له دې حقيقت څخه انکار نۀ شي کېدای چې رياضيپوهان د ثبوت پر وخت \emph{ډېر ځله} انځورونه کاروي۔ (په تقريبي ټکو: ښه رياضيپوهان پوهېږي چې کله پر هندسي وجدان تکيه کولے شي۔)

په لنډه توګه: پر تشې شورماشور باور مۀ کوئ؛ يا لږ تر لږه يې بې له پوښتنې مۀ منئ۔ د دې په اړه د نور لوست لپاره \citet{Giaquinto2007} کتلے شئ۔

\end{document}
